\begin{frame}{Patterns worth noticing}
  \begin{columns}[T,onlytextwidth]
    \begin{column}{0.31\textwidth}
      \begin{block}{Zeros}
        sparse, triangular, diagonal
      \end{block}
      \begin{block}{Repetition}
        equal rows, columns, blocks
      \end{block}
    \end{column}
    \begin{column}{0.31\textwidth}
      \begin{exampleblock}{Dependence}
        proportional rows or columns
      \end{exampleblock}
      \begin{exampleblock}{Symmetry}
        mirrored entries
      \end{exampleblock}
    \end{column}
    \begin{column}{0.31\textwidth}
      \begin{alertblock}{Use the pattern}
        State what becomes simpler and why.
      \end{alertblock}
    \end{column}
  \end{columns}
\end{frame}

\begin{frame}{Checkpoint: see before you compute}
  \begin{center}
  \begin{tikzpicture}[node distance=0.65cm]
    \tikzset{c/.style={rounded corners=5pt,draw=line,very thick,fill=paper,minimum width=2.2cm,minimum height=0.9cm,align=center,font=\bfseries}}
    \node[c] (look) {look for\\patterns};
    \node[c,right=of look] (explain) {explain their\\meaning};
    \node[c,right=of explain] (use) {use them to\\simplify};
    \draw[-{Latex[length=3mm]},thick,blue] (look)--(explain);
    \draw[-{Latex[length=3mm]},thick,blue] (explain)--(use);
  \end{tikzpicture}
  \end{center}
  \vspace{0.7cm}
  \concept{Noticing structure is mathematical understanding.}
\end{frame}
